\chapter{Feature Importance}
\label{ch:featureimportance}

Having established how the models are trained and evaluated, this chapter turns
to the interpretive question at the heart of the thesis: which financial features
do the fitted models rely on? Three importance methods are applied to each of the
three models --- native importance, permutation importance, and SHAP. They are
introduced here as methods; their numerical results, and the divergence between
them, are reported in Chapter~\ref{ch:results}.

One caveat frames everything that follows and is repeated deliberately. Feature
importance is a statement about the model, not about the economy. It measures
what a trained classifier uses to separate bankrupt from healthy firms; it does
not identify the causal mechanisms that drive firms into insolvency. Without a
design for causal identification, every ranking in this and the following chapter
is to be read as \emph{associative} --- informative about predictive reliance,
silent about cause. This distinction is not a rhetorical hedge but a genuine
limit of the exercise, and it becomes especially important where the three
methods disagree.

\section{Native Importance}
\label{sec:fi-native}

Native importance reads the importance directly off the fitted model, and it
therefore means something different for each model family. For the logistic
regression, the natural measure is the vector of estimated coefficients from
Equation~\eqref{eq:logit}. Because the features are normalised to a common scale,
the magnitude of a coefficient reflects the strength of its influence on the
predicted log-odds, and its sign reflects direction: a positive coefficient
pushes a firm towards the bankruptcy class, a negative one away from it. Signed,
directional importance is a distinctive strength of the linear model.

For the two tree ensembles, native importance is the Gini impurity importance:
each feature is credited with the total reduction in node impurity it produces,
summed over all splits and averaged over all trees. This measure is always
non-negative and hence \emph{unsigned}; it records that a feature matters, not in
which direction. It also carries a well-documented bias. \textcite{strobl2007}
show that impurity importance systematically favours continuous and
high-cardinality features, because such features offer more candidate split
points and can therefore reduce impurity by chance more often than coarse or
binary features. On a dataset of $95$ continuous ratios of differing granularity,
this bias is a live concern, and it is one reason the analysis does not stop at
native importance.

\section{Permutation Importance}
\label{sec:fi-permutation}

Permutation importance takes a model-agnostic, performance-based view. The values
of a single feature are randomly shuffled across observations, breaking its
association with the target while leaving all other features intact, and the
resulting drop in the $F_2$-score is recorded. A feature the model genuinely
relies on will, when scrambled, cause a large deterioration; an irrelevant
feature will cause none. To average out the randomness of the shuffle, the
procedure is repeated ten times and the mean drop, together with its standard
deviation, is reported. Because it is defined purely in terms of the loss, the
method applies identically to all three models and thus provides a common
comparison anchor.

Its principal weakness is also well understood and is directly relevant here.
When features are strongly correlated, permuting one of them barely hurts the
model, because a correlated partner still carries almost the same information and
substitutes for the scrambled feature. On the highly collinear Taiwanese ratios,
this ``correlation masking'' can drive the measured importances towards
implausibly small values, so that even features the model demonstrably uses
appear unimportant. This behaviour is not a defect of the implementation but an
intrinsic property of permutation under dependence, and it motivates a third
method whose theoretical guarantees are stronger.

\section{SHAP}
\label{sec:fi-shap}

SHAP (SHapley Additive exPlanations) grounds feature importance in cooperative
game theory and, unlike the previous two methods, comes with axiomatic
guarantees \parencite{lundberg2017}. Its starting point is the idea of an
\emph{explanation model}: a simple, interpretable surrogate that approximates the
complex model locally. \textcite{lundberg2017} define the class of additive
feature-attribution methods as those whose explanation model is a linear function
of binary indicators,
\begin{equation}
  g(z') = \phi_0 + \sum_{i=1}^{M} \phi_i z_i',
  \label{eq:shap-additive}
\end{equation}
where $z' \in \{0,1\}^M$ indicates the presence or absence of each of the $M$
simplified features, $\phi_0$ is the model output with no feature information,
and each $\phi_i$ is the contribution attributed to feature $i$. The prediction
is thus decomposed additively into a base value plus one term per feature.

The question is how to assign the attributions $\phi_i$. SHAP answers it with the
Shapley value, the unique allocation from cooperative game theory that
distributes the ``payout'' of a prediction fairly among the features. Treating
the features as players in a game whose payoff is the model output, the
contribution of feature $i$ is its marginal effect averaged over every coalition
of the remaining features. Writing $F$ for the full set of features, a model
$f_{S \cup \{i\}}$ trained with feature $i$ present is compared against a model
$f_S$ trained without it, over all subsets $S \subseteq F \setminus \{i\}$:
\begin{equation}
  \phi_i = \sum_{S \subseteq F \setminus \{i\}}
  \frac{|S|!\,\bigl(|F| - |S| - 1\bigr)!}{|F|!}
  \left[\, f_{S \cup \{i\}}\!\left(x_{S \cup \{i\}}\right) - f_S\!\left(x_S\right) \right].
  \label{eq:shap-shapley}
\end{equation}
The combinatorial weight in Equation~\eqref{eq:shap-shapley} counts the
orderings in which feature $i$ joins coalition $S$, so that $\phi_i$ is a weighted
average of the feature's marginal contribution across all possible coalitions.
This is what gives SHAP its fairness properties, but it is also what makes the
exact computation intractable: the sum ranges over all subsets of the remaining
features, of which there are $2^{|F|-1}$. With $95$ features this amounts to
$2^{94}$ coalitions per feature --- a number far beyond any feasible computation.

SHAP is therefore made practical by exploiting model structure to evaluate
Equation~\eqref{eq:shap-shapley} exactly without enumerating the coalitions. For
the logistic regression, a linear explainer computes the Shapley values in closed
form from the coefficients and the feature distribution. For the two tree
ensembles, a tree explainer traverses the trees to obtain the exact Shapley
values in polynomial rather than exponential time. In both cases the values are
computed for the bankruptcy class over all firms, and a feature's overall
importance is summarised by the mean absolute SHAP value across firms.

Two properties make SHAP the interpretive centre of gravity for this thesis.
First, it returns \emph{signed} attributions for \emph{all three} models: the
logistic coefficients are signed but only linear, and the impurity importance is
unsigned, whereas SHAP places every model on the same signed, additive axis and
so makes their rankings directly comparable. Second, because it accounts for
feature coalitions rather than perturbing features one at a time, it does not
collapse under the correlation that so weakens permutation importance. Where the
two performance-based views disagree, as they will in Chapter~\ref{ch:results},
SHAP provides the more trustworthy reading, while permutation importance remains
valuable precisely because its collapse diagnoses the collinearity of the data.
